\documentclass[11pt,letterpaper]{article}
\usepackage[margin=1.05in]{geometry}
\usepackage{amsmath,amssymb,amsthm,xcolor}
\usepackage[colorlinks=true,linkcolor=teal!60!black,citecolor=teal!60!black,urlcolor=teal!60!black]{hyperref}
\hypersetup{pdftitle={LogCFL is unambiguous: Hoza's isolation for proof trees}}
\newtheorem{theorem}{Theorem}
\newtheorem{lemma}[theorem]{Lemma}
\newtheorem{corollary}[theorem]{Corollary}
\theoremstyle{definition}
\newtheorem{definition}[theorem]{Definition}
\theoremstyle{remark}
\newtheorem{remark}[theorem]{Remark}
\newcommand{\N}{\mathbb N}
\newcommand{\LogCFL}{\mathrm{LogCFL}}
\newcommand{\UAux}{\mathrm{UAuxPDA}}
\newcommand{\Cert}{\mathrm{Cert}}
\setlength{\parindent}{0pt}
\setlength{\parskip}{5pt plus 1pt minus 1pt}
\raggedbottom
\emergencystretch=3em
\title{\bfseries LogCFL is unambiguous:\\Hoza's isolation for proof trees}
\author{A research note extending Hoza's proof of NL = UL}
\date{November 2026}
\begin{document}
\maketitle
\begin{abstract}
Hoza recently proved $\mathrm{NL}=\mathrm{UL}$ by giving a logspace generator
with an $O(\log n)$-bit seed whose weights make the minimum-weight path between
every pair of vertices of a layered digraph unique, for some seed. We observe
that the same construction isolates minimum-weight \emph{proof trees} in
layered semi-unbounded circuits of logarithmic depth. The 2-adic valuation of the layer
index in Hoza's weights is replaced by the AND-level of the gate, and the
analysis goes through almost verbatim. Combined with the unambiguous
certificate machines of van Melkebeek and Prakriya, this gives
\[
 \LogCFL \;=\; \UAux(\log n,\,n^{O(1)}) \;=\; \mathrm{co}\UAux(\log n,\,n^{O(1)}).
\]
Previously this was known with $O(\log^{3/2} n)$ space, or with polynomial advice.
\end{abstract}

\section{Introduction}
$\LogCFL$ is the class of languages logspace-reducible to context-free
languages. It equals $\mathrm{AuxPDA}(\log n,n^{O(1)})$, the languages accepted
by nondeterministic machines with $O(\log n)$ work space, an unbounded stack
and polynomial running time \cite{Sudborough}. It also equals
logspace-uniform $\mathrm{SAC}^1$ \cite{Venkateswaran}. $\UAux(\log
n,n^{O(1)})$ is the subclass accepted by such machines that have at most one
accepting computation on every input.

Reinhardt and Allender combined the isolation lemma with double inductive
counting and proved $\LogCFL/\mathrm{poly}=\UAux/\mathrm{poly}$, alongside
$\mathrm{NL}/\mathrm{poly}=\mathrm{UL}/\mathrm{poly}$ \cite{RA}. Van Melkebeek and Prakriya derandomized
the isolation step partially and obtained $\LogCFL\subseteq\UAux$ in space
$O(\log^{3/2}n)$ \cite{vMP}. Hoza's weight assignment generator \cite{Hoza}
removes the remaining loss for layered digraphs and proves $\mathrm{NL}=\mathrm{UL}$.

\begin{theorem}\label{thm:main}
$\LogCFL = \UAux(\log n,n^{O(1)}) = \mathrm{co}\UAux(\log n,n^{O(1)})$.
\end{theorem}

The new ingredient is Theorem~\ref{thm:generator}, an explicit
weight assignment generator for proof trees. Its proof follows
Hoza~\cite[\S2]{Hoza} closely. In a layered digraph, Hoza uses the
2-adic valuation of the layer index to single out one dominant internal
layer of every interval, and to ensure that a path meets at most $2^\Delta$ layers of
valuation $M-\Delta$. A proof tree of a circuit of logarithmic depth has
this structure natively: below an AND gate at AND-level $k$ it contains at
most $2^m$ copies of AND gates at AND-level $k-m$. So we can scale the
weights by the AND-level itself. The rest of the argument, including the
constant $6$, carries over. Section~\ref{sec:machines} then plugs the
generator into the machinery of \cite{vMP}, in the same way that Hoza plugs
his into that of \cite{RA}.

\section{Circuits and proof trees}

\begin{definition}[Normal form]\label{def:normal}
A \emph{normal-form circuit} of AND-depth $K$ has layers $V_0,V_1,\dots,V_{2K}$, each
of size at most $w$; the gates of $V_i$ are indexed by $[w]$.
\begin{itemize}
\item Layer $V_0$ consists of leaves labeled by literals or constants.
\item Each odd layer $V_{2k-1}$ consists of AND gates of fan-in $1$ or $2$, with
 inputs in $V_{2k-2}$. Both inputs may be the same gate.
\item Each even layer $V_{2k}$ with $k\ge 1$ consists of OR gates of arbitrary
 positive fan-in, with inputs in $V_{2k-1}$.
\end{itemize}
An AND gate in $V_{2k-1}$ has \emph{AND-level} $k$. Let $n$ denote the
number of gates.
\end{definition}

The important feature is that the layer next to the leaves consists of AND
gates, so every OR gate has only AND gates as inputs. Without it,
isolation fails: in $x_1\vee x_2$ with both inputs true, the gate has two proof
trees, each of weight zero, whatever the weights of the gates are.

\begin{definition}[Proof trees]\label{def:cert}
Fix an input $z$. For a gate $g$ define a finite set $\Cert(g)$ of labeled
trees inductively. The tree $\langle g\rangle$ stands for $g$ alone, and
$\langle g;F_1,\dots\rangle$ for $g$ with subtrees $F_1,\dots$.
\begin{itemize}
\item A leaf that is true on $z$ has $\Cert(g)=\{\langle g\rangle\}$. A false
 leaf has $\Cert(g)=\emptyset$.
\item An AND gate with inputs $a,b$ has $\Cert(g)=\{\langle g;F,F'\rangle: F\in\Cert(a),\,F'\in\Cert(b)\}$.
 This applies also when $a=b$, in which case $F$ and $F'$ are chosen independently.
 With a single input $a$, $\Cert(g)=\{\langle g;F\rangle:F\in\Cert(a)\}$.
\item An OR gate with inputs $h_1,\dots,h_r$ has $\Cert(g)=\bigcup_i \{\langle g;F\rangle: F\in\Cert(h_i)\}$.
\end{itemize}
Gate $g$ evaluates to $1$ on $z$ if and only if $\Cert(g)\neq\emptyset$.
Given a weight $\rho\colon V\to\N$, the weight $\rho(F)$ of a tree is the sum
of $\rho$ over all of its nodes, with multiplicity. Let
$\mu^\rho(g)=\min_{F\in\Cert(g)}\rho(F)$. We say that $\rho$ is
\emph{min-isolating} for $(C,z)$ if every gate $g$ with $\Cert(g)\ne\emptyset$
has exactly one $F\in\Cert(g)$ with $\rho(F)=\mu^\rho(g)$.
\end{definition}

This agrees with the notion of certificates and their weights in
\cite[Definition~6]{vMP}. Proof trees of a gate are trees over the unfolding
of the circuit, so a gate may occur in one proof tree several times.

\begin{lemma}[Level profile]\label{lem:profile}
Let $u$ be an AND gate at AND-level $k$ and $F\in\Cert(u)$. For every $m\ge1$, $F$
contains at most $2^m$ nodes labeled by AND gates of AND-level $k-m$, and it contains exactly one
node of AND-level $k$, namely its root.
\end{lemma}
\begin{proof}
Below an AND node at layer $2k'-1$ there are at most two OR nodes at
layer $2k'-2$, if $k'\ge2$, and each OR node has exactly one child, an AND node at layer
$2k'-3$. Induction on $m$.
\end{proof}

\section{The generator}

Let $\mathcal C\colon[w]\to\{0,1\}^\ell$ be a code with minimum distance $d$,
computable in space $O(\log(wK))$, with $\ell=O(\log(wK))$. Justesen codes
\cite{Justesen} achieve $d=\Theta(\log(wK))$ with these parameters, as used in
\cite{Hoza}. For a gate $g$ write $\mathcal C(g)$ for the codeword of its index
within its layer.

\begin{definition}[Weights]\label{def:weights}
For a seed $x\in\{0,1\}^\ell$, let an AND gate $g$ at AND-level $k$ have weight
\[
 \rho_x(g)=\sum_{j=1}^{\ell} 6^{\,k+\ell-j}\cdot\mathcal C(g)_j\cdot x_j ,
\]
and let OR gates and leaves have weight $0$.
\end{definition}

This is Hoza's Eq.~(1) \cite{Hoza}, with $\nu_2(i)$ replaced by $k$. Weights are at
most $6^{K+\ell}$, and by Lemma~\ref{lem:profile} every proof tree weighs less than
$6^{K+\ell}$. For $K,\ell=O(\log n)$ this is $n^{O(1)}$. Given $x$ and $g$, the
weight is computable in space $O(\log n)$.

\begin{theorem}\label{thm:generator}
Let $C$ be a normal-form circuit with $n$ gates, and let $z$ be an input. If $x$ is
uniformly random, then
\[
 \Pr_x[\rho_x \text{ is not min-isolating for }(C,z)]\ \le\ 2^{-d}\cdot n\,w^2 .
\]
In particular, if $d\ge 3\log_2 n+1$, then some seed is min-isolating.
\end{theorem}

Fix an AND gate $u$ at AND-level $k$ with $\Cert(u)\ne\emptyset$, and put
$B_j=6^{k+\ell-j}$. For $f\colon\{0,1\}^\ell\to\mathbb R$ write $\partial_j
f(x)=f(x^{j\leftarrow1})-f(x^{j\leftarrow0})$.

\begin{lemma}[Derivatives; cf.\ {\cite[Lemma~2.3]{Hoza}}]\label{lem:deriv}
For every $j\in[\ell]$ and every seed $x$,
\[
 \partial_j\mu^{\rho_x}(u)\in\bigl[B_j\,\mathcal C(u)_j,\ B_j\,(\mathcal C(u)_j+\tfrac12)\bigr].
\]
\end{lemma}
\begin{proof}
First fix $F\in\Cert(u)$ and let $f_F(x)=\rho_x(F)$. Flipping $x_j$ from $0$ to
$1$ increases the root's weight by exactly $B_j\mathcal C(u)_j$. It increases
the weight of each AND node at level $k-m$ by at most $6^{-m}B_j$, and by
Lemma~\ref{lem:profile} there are at most $2^m$ such nodes. All increases are
nonnegative. Hence
\[
 \partial_j f_F(x)\in\Bigl[B_j\mathcal C(u)_j,\ B_j\mathcal C(u)_j+B_j\sum_{m\ge1}2^m6^{-m}\Bigr]
 =\bigl[B_j\mathcal C(u)_j,\ B_j(\mathcal C(u)_j+\tfrac12)\bigr].
\]
Now let $F_0$ and $F_1$ be minimizers for the seeds $x^{j\leftarrow0}$ and
$x^{j\leftarrow1}$. Then
\[
 \partial_j\mu(u)=\bigl(f_{F_1}(x^{j\leftarrow1})-f_{F_0}(x^{j\leftarrow1})\bigr)+\partial_jf_{F_0}(x)\le\partial_jf_{F_0}(x)
\]
and
\[
 \partial_j\mu(u)=\partial_jf_{F_1}(x)+\bigl(f_{F_1}(x^{j\leftarrow0})-f_{F_0}(x^{j\leftarrow0})\bigr)\ge\partial_jf_{F_1}(x),
\]
since the first bracket is at most $0$ and the second at least $0$. Both bounds then lie in the interval.
\end{proof}

The exchange argument uses only that $\mu$ is a pointwise minimum of functions with the
stated derivative bounds. It does not depend on the objects being paths or trees.

\begin{lemma}[Collisions; cf.\ {\cite[Theorem~2.4]{Hoza}}]\label{lem:collide}
Let $u\ne v$ be AND gates in the same layer with $\Cert(u),\Cert(v)\neq\emptyset$, and let
$A=\{j:\mathcal C(u)_j=\mathcal C(v)_j\}$. Then the set
$E=\{x:\mu^{\rho_x}(u)=\mu^{\rho_x}(v)\}$ has $|E|\le 2^{|A|}\le 2^{\ell-d}$.
\end{lemma}
\begin{proof}
Both gates have the same AND-level, so they share the scales $B_j$. Let $x\in
E$, and let $y\ne x$ agree with $x$ on $A$. Let $j^*$ be the first coordinate in
which $x$ and $y$ differ, so $j^*\notin A$. Flip the differing coordinates one at a
time in increasing order. By Lemma~\ref{lem:deriv}, the first flip changes the
value of the gate with codeword bit $1$ at $j^*$ by at least $B_{j^*}$, and the other by at most
$B_{j^*}/2$, both in the same direction. So afterwards $|\mu(u)-\mu(v)|\ge
B_{j^*}/2$. Every later flip at a coordinate $j>j^*$ moves both values in the same direction
by amounts in $[0,\frac32B_j]$, so it changes the difference by at most $\frac32B_j$. In
total, these flips change the difference by less than
$\frac32B_{j^*}\sum_{t\ge1}6^{-t}=\frac3{10}B_{j^*}<\frac12B_{j^*}$. Hence
$y\notin E$. So distinct elements of $E$ differ somewhere in $A$, and $|E|\le
2^{|A|}$.
\end{proof}

\begin{proof}[Proof of Theorem~\ref{thm:generator}]
Call $x$ \emph{bad} if some OR gate $g$ has two distinct inputs $u,v$ with
$\Cert(u),\Cert(v)\ne\emptyset$ and $\mu(u)=\mu(v)$. The inputs of an OR gate
are AND gates of one layer, so Lemma~\ref{lem:collide} and a union bound over at most $n$ OR
gates and $\binom w2$ pairs bound the probability of a bad seed by $2^{-d}nw^2/2$.

Let $x$ be good. We show by induction on the layer that every gate $g$ with
$\Cert(g)\ne\emptyset$ has a unique minimizer.
\begin{itemize}
\item For a true leaf, this holds trivially.
\item For an AND gate, $\rho(\langle g;F,F'\rangle)=\rho(g)+\rho(F)+\rho(F')$
 with independent choices, so the minimizers are exactly the pairs of minimizers,
 and there is one.
\item For an OR gate, $\rho(\langle g;F\rangle)=\rho(F)$. By goodness, exactly
 one satisfied input $h$ attains $\min_h\mu(h)$, and $h$ has a unique minimizer. \qedhere
\end{itemize}
\end{proof}

\begin{remark}
The base $6$ is the smallest integer for which both steps work: the derivative bound needs
$\sum_m(2/b)^m\le\frac12$, that is, $b\ge6$, and the flip step needs
$\frac12>\frac32\cdot\frac1{b-1}$, that is, $b>4$. Exhaustive checks on small circuits
find violations of Lemma~\ref{lem:deriv} for bases $2$ to $5$; see the accompanying
script.
\end{remark}

\section{Unambiguous machines}\label{sec:machines}

\begin{proof}[Proof of Theorem~\ref{thm:main}]
The inclusion $\UAux\subseteq\mathrm{AuxPDA}=\LogCFL$ is trivial, so let $L\in\LogCFL$.
\begin{enumerate}
\item \emph{Reduction.} By \cite[Proposition~2]{vMP}, a logspace mapping reduction sends an
 input of $L$ to a layered alternating semi-unbounded circuit of logarithmic depth
 and an input $z$, whose output gate evaluates to $1$ exactly when the original
 input is in $L$. The transformation in its proof (\cite[Claim~6]{vMP}) produces
 exactly the normal form of Definition~\ref{def:normal}, padded with fan-in-1 AND
 gates where necessary \cite[footnotes~10 and~11]{vMP}. In particular $K=O(\log n)$,
 and Theorem~\ref{thm:generator} applies with $\ell=O(\log n)$ and $d\ge3\log_2n+1$.
\item \emph{Certificate machines.} By \cite[Lemma~3]{vMP}, there are unambiguous
 machines $\mathrm{WeightCheck}$ and $\mathrm{WeightEval}$, each with a stack not charged to
 the space bound, such that for every layered semi-unbounded circuit $C$ of depth
 $D$ with $n$ gates, input $z$, gate weights $\rho$ and gate $g$:
 $\mathrm{WeightCheck}(C,z,\rho)$ decides whether $\rho$ is min-isolating for
 $(C,z)$, and $\mathrm{WeightEval}(C,z,\rho,g)$ computes $\mu^\rho(g)$, with
 the value $\infty$ when $g$ is false, provided $\rho$ is min-isolating. Both run in time
 $\mathrm{poly}(2^D,\log W,n)$ and space $O(D+\log W+\log n)$. Here $D=2K=O(\log n)$ and
 $W<6^{K+\ell}=n^{O(1)}$. The weights $\rho_x$ are computed on demand in logspace.
\item \emph{Seed search.} For the seeds $x\in\{0,1\}^\ell$ in lexicographic order, run
 $\mathrm{WeightCheck}(C,z,\rho_x)$. Stop at the first $x$ it accepts, which exists by
 Theorem~\ref{thm:generator}. Then accept if and only if
 $\mathrm{WeightEval}(C,z,\rho_x,\mathrm{out})<\infty$.
\end{enumerate}
There are $2^\ell=n^{O(1)}$ seeds, so the total time is polynomial and the space is
$O(\log n)$. Every call is unambiguous in the sense of \cite{vMP,RA}: exactly one
computation path returns, and it returns the correct answer, while all other paths reject.
The calls run sequentially with the stack empty between them. So the
composed machine has exactly one surviving path, and that path outputs the
correct answer. This gives $L\in\UAux\cap\mathrm{co}\UAux$.
\end{proof}

This is the same scheme as Hoza's application of \cite{RA} to $\mathrm{NL}$. The
prefix-by-prefix seed construction of \cite[Algorithm~1]{vMP} is unnecessary here,
because the whole seed has $O(\log n)$ bits.

\section{Remarks}

\paragraph{Consequences.} With count-preserving translations between AuxPDAs
and $\mathrm{SAC}^1$ circuits \cite{Vinay,vMP}, Theorem~\ref{thm:main} says that
every language in logspace-uniform $\mathrm{SAC}^1$ has such a circuit whose output gate
has at most one proof tree. As with $\mathrm{NL}\subseteq\oplus\mathrm{L}$,
it follows that $\LogCFL\subseteq\oplus\mathrm{SAC}^1$ and
$\LogCFL\subseteq\mathrm{Mod}_p\mathrm{SAC}^1$, uniformly. Reinhardt and
Allender list the question whether $\LogCFL$ reduces to an
\emph{unambiguous context-free language} \cite[\S4]{RA}. That stronger form of
unambiguity is not addressed here.

\paragraph{Novelty.} The argument is a direct adaptation and needs no new technique.
To our knowledge it has not been stated before. We have not seen any revision of
\cite{Hoza} or follow-up that mentions $\LogCFL$.

\paragraph{Verification and disclosure.} This note was written with AI assistance at
the request of Christian Reitwiessner. The proof was found by an AI agent and
checked step by step by a second, independent AI review. That review found that the
AND-first-layer condition is necessary, and that the level profile is an upper bound
rather than an equality. Both corrections are incorporated here. Exhaustive seed
enumeration on small adversarial circuits confirms Lemmas~\ref{lem:deriv}
and~\ref{lem:collide}; see \texttt{check.py}. The appendix machinery of
\cite{vMP} (Lemma~3) and the logspace computability of Justesen codes are used as
published. \textbf{The note has not been reviewed by a human expert.}

\begin{thebibliography}{9}
\bibitem{Hoza} W.~M.~Hoza. NL = UL. ECCC TR26-323, 2026. \url{https://eccc.weizmann.ac.il/report/2026/323/}
\bibitem{Justesen} J.~Justesen. A class of constructive asymptotically good algebraic codes. \emph{IEEE Trans.\ Inform.\ Theory} 18(5):652--656, 1972.
\bibitem{RA} K.~Reinhardt and E.~Allender. Making nondeterminism unambiguous. \emph{SIAM J.\ Comput.} 29(4):1118--1131, 2000. \url{https://people.cs.rutgers.edu/~allender/papers/nlul.pdf}
\bibitem{Sudborough} I.~H.~Sudborough. On the tape complexity of deterministic context-free languages. \emph{J.\ ACM} 25(3):405--414, 1978.
\bibitem{vMP} D.~van Melkebeek and G.~Prakriya. Derandomizing isolation in space-bounded settings. \emph{SIAM J.\ Comput.} 48(3):979--1021, 2019; ECCC TR17-052. \url{https://pages.cs.wisc.edu/~dieter/Papers/r-ul-eccc.pdf}
\bibitem{Venkateswaran} H.~Venkateswaran. Properties that characterize LOGCFL. \emph{J.\ Comput.\ System Sci.} 43(2):380--404, 1991.
\bibitem{Vinay} V.~Vinay. Counting auxiliary pushdown automata and semi-unbounded arithmetic circuits. In \emph{Proc.\ 6th Structure in Complexity Theory Conference}, 270--284, 1991.
\end{thebibliography}
\end{document}
